%قبل از شروع مسئله ترجیحاً چیزی ننویسید. اما اگر خواستید بنویسید، در قالب لیستی (فقرات و شمارش) بنویسید.
%برای اضافه کردن نمره، بعدی از شروع‌مسئله، [x] را اضافه کنید و به جای x  نمره مورد نظر را قرار دهید. برای زیرمسئله‌ها هم به همین صورت است.
%\شروع‌مسئله[10]
\شروع‌مسئله [18]
%\متن‌سیاه{سوال عملی اول:}\\
%\متن‌سیاه{دراپ‌تست}\\
محمد که در حال مرور درس های گذشته‌اش بود، با یک سوال مواجه می‌شود. صورت سوال به شکل زیر است:

مسئله مرتب سازی را در نظر بگیرید. فرض کنید اعداد به صورت دقیق مشخص نیستند. در عوض برای هر عدد، یک بازه‌ی بسته روی اعداد حقیقی داریم که می‌دانیم آن عدد متعلق به آن بازه است. یعنی ‌\متن‌لاتین{n} بازه به صورت 
$[a_i, b_i]$
داریم، که در آن 
$a_i \le b_i$. 
هدف این است که این بازه ها مرتب کنیم؛ یعنی یک جایگشت 
$\langle i_1, i_2, ..., i_n \rangle$
از بازه‌ها بسازیم به طوری که برای هر 
$j = 1, 2, ..., n$
مقادیر
$c_i \in [a_{i_j}, b_{i_j}]$
موجود باشند که در نامساوی
$c_1 \le c_2 \le ... \le c_n$
صدق کنند.
فرض کنید که هر بازه، حداقل با 
$d-1$
بازه‌ی دیگر هم‌پوشانی داشته باشد. الگوریتمی از 
$O(n\log(\frac{n}{d}))$
ارائه دهید که این \متن‌لاتین{n} بازه با درجه‌ی \متن‌لاتین{d} را مرتب کند. بدین ترتیب، اگر
$d = \Theta(n)$
باشد، الگوریتم در زمان
$O(n)$
اجرا خواهد شد. هم‌چنین اگر
$d = \Theta(1)$
باشد، الگوریتم در زمان
$O(n\log(n))$
اجرا خواهد شد.

حال محمد که دیگر کلا این چیز ها یادش رفته است، از شما می‌خواهد که راه این سوال را به او بگویید. اگر هم محمد اشتباه می‌کند باید به او ثابت کنید که این کار شدنی نیست. 
\solution{
پاسخ. این کار شدنی نیست. برای اثبات این کار، از برهان خلف استفاده می‌کنیم. فرض کنید که روشی با شرایط گفته شده وجود داشته باشد. می‌خواهیم یک الگوریتم مرتب‌سازی برای n عدد ارائه دهیم که از 
$ O(n) $
باشد. فرض کنید که n عدد به شکل 
$A_1, A_2, ..., A_n$
داریم. مینیمم این اعداد را
$A_x$
و ماکسیمم این اعداد را 
$A_y$
می‌نامیم. حال n بازه متناظر این اعداد درست می‌کنیم که بازه‌ی i ام آنها به شکل 
$a_i = b_i = A_i$
تعریف می‌کنیم. به n بازه‌ای که ساختیم، n بازه‌ی دیگر اضافه می‌کنیم که بازه‌ی n+i ام بدین شکل است:
$a_{n+i}=A_x - 1, b_{n+i} = A_y+1$.
حال این 2n بازه را درنظر بگیرید. می‌دانیم d تعریف شده در الگوریتم حداقل برابر 
$n$ است. 
طبق الگوریتم داده شده، می‌توانیم در اردر زیر این بازه‌ها را مرتب کنیم.
$$T = O(2n\log(\frac{2n}{d})) = O(n\log(2)) = O(n)$$
حال کافی است از جایگشت مرتب‌شده‌ی بازه‌ها، بازه‌های با شماره های 
$n+1$
تا
$2n$
را حذف کنیم. بازه‌های باقی مانده، همان اعدادی هستند که می‌خواستیم آن‌ها را مرتب کنیم. پس توانستیم 
$n$
عدد را در 
$O(n)$
مرتب کنیم. می‌دانیم این کار امکان‌پذیر نیست و مرتب‌سازی مقایسه‌ای باید از 
$\Omega(n\log(n))$
باشد.

}
%\متن‌سیاه{ورودی:}\\
%در خط اول ورودی، ابتدا k و سپس m داده می‌شود.\\
%\متن‌سیاه{خروجی:}\\
%در تنها سطر خروجی، تعداد آیفون‌های مورد نیاز می‌آید. اگر در هیچ صورتی چنین کاری ممکن نبود ۱- چاپ می‌شود.\\
%\متن‌سیاه{محدودیت‌ها:}\\
%$1\leq k\leq 10^4$

%$m\leq10^9$
%\\
%\textbf{\متن‌سیاه{ورودی و خروجی نمونه}\\
%% تست اول
%\begin{latin}
%	\begin{tabularx}{\textwidth}{|X | X|}
%		\hline
%		\textbf{stdin} &
%		\textbf{stdout}\\
%		\hline
%		100 100&
%		1\\
%		\hline
%	
%	\end{tabularx}
%\end{latin}
%
%% تست دوم
%
%\begin{latin}
%	\begin{tabularx}{\textwidth}{|X | X|}
%		\hline
%		\textbf{stdin} &
%		\textbf{stdout}\\
%		\hline
%		340 20&
%		3\\
%		\hline
%		
%	\end{tabularx}
%\end{latin}}
% صورت اصلی مسئله را در این‌جا بنویسید.
%اگر مسئله چند قسمت داشت به این صورت عمل کنید:
%\شروع{زیرمسائل}
%\زیرمسئله[0] % مسئلۀ ۱آ
%\زیرمسئله % مسئلۀ ۱ب
%\زیرمسئله % مسئلۀ ۱پ
%\زیرمسئله % مسئلۀ ۱ت
%\پایان{زیرمسائل}




%%%%%%%%%%%%%%%%%%%%%%%%%%%%%%%%%%%%%%%%%%%%%%%%%%%%%
% برخی دستورات مفید برای تمرین‌ها. %
%%%%%%%%%%%%%%%%%%%%%%%%%%%%%%%%%%%%%%%%%%%%%%%%%%%%%

% فرمول‌های نمونه:
%$f(x) = \Theta(x)$
%$T(x, y) \leq \log(x) + 2^y + \binom{2n}{n}$
% {\tt code\_function}


% لیست موارد بدون شماره‌گذاری:
%\شروع{فقرات}
%    \فقره مورد اول 
%        \شروع{فقرات}
%            \فقره مورد اول داخلی 
%                \شروع{فقرات}
%                    \فقره اولین مورد خیلی داخلی 
%                    \فقره دومین مورد خیلی داخلی 
%                \پایان{فقرات}
%            \فقره مورد دوم داخلی 
%        \پایان{فقرات}
%    \فقره مورد دوم 
%\پایان{فقرات}

% لیست موارد شماره‌گذاری شده:
%\شروع{شمارش}
%    \فقره یک مورد 
%    \فقره یک مورد دیگر 
%    \فقره باز هم یک مورد دیگر
%\پایان{شمارش}

% برای نوشتن یک فرمول چند خطی که یک شماره بخوره. برای حالت‌های دیگر به سایت زیر مراجعه کنید:
% https://www.overleaf.com/learn/latex/Aligning_equations_with_amsmath
%\begin{align} 
%    \begin{split}
%        (x+y)^3 &= (x+y)^2(x+y) \\
%                &= (x^2+2xy+y^2)(x+y) \\
%                &= (x^3+2x^2y+xy^2) + (x^2y+2xy^2+y^3) \\
%                &= x^3+3x^2y+3xy^2+y^3
%    \end{split}                                 
%\end{align}

% نوشتن ماتریس و آرایه دوبعدی
%\begin{align}
%    A = 
%    \begin{bmatrix}
%        A_{11} & A_{21} \\
%        A_{21} & A_{22}
%    \end{bmatrix}
%\end{align}

% اضافه کردن تصویر با شرح. برای اضافه کردن تصویر بدون شرح، فقط دستور درج‌تصویر را استفاده کنید.
%\شروع{شکل}[h!]
%\تنظیم‌ازوسط
%  \درج‌تصویر[پهنا=0.5\پهنای‌متن]{img.jpg}
%\شرح{شرح تصویر}
%\برچسب{تصویر:۱}
%\پایان{شکل}
%\پایان{زیرمسائل}

%کد زیر را ببینید.
%%\lstinputlisting{source_filename.py}
%\begin{latin}
%\begin{lstlisting}
%# Import the modules
%import sys
%import random
%
%ans = True
%
%while ans:
%question = raw_input("Ask the magic 8 ball a question: (press enter to quit) ")
%
%answers = random.randint(1,8)
%
%if question == "":
%sys.exit()
%
%elif answers == 1:
%print "It is certain"
%
%elif answers == 2:
%print "Outlook good"
%
%elif answers == 3:
%print "You may rely on it"
%
%elif answers == 4:
%print "Ask again later"
%
%elif answers == 5:
%print "Concentrate and ask again"
%
%elif answers == 6:
%print "Reply hazy, try again"
%
%elif answers == 7:
%print "My reply is no"
%
%elif answers == 8:
%print "My sources say no"
%\end{lstlisting}
%\end{latin}
